
\subsection{Experiment Timeline and Data}

Below, we detail the timeline of the experiment and main data collection activities. 

\begin{figure}[H]
\centering
\begin{tikzpicture}[
  x=1cm, y=1cm, >=Latex,
  event/.style={draw, rounded corners=2pt, align=center, inner sep=5pt,
                minimum width=3.25cm, minimum height=1.05cm},
  arrow/.style={-{Latex[length=2.2mm]}, thick}
]
  \node[event, fill=gray!8] (census) at (0,0)
    {Household census\\and baseline survey};
  \node[event, fill=blue!7] (outreach) at (4.1,0)
    {Community outreach\\and SMS enrollment};
  \node[event, fill=green!8] (campaign) at (8.2,0)
    {Deworming campaign\\administrative take-up};
  \node[event, fill=orange!10] (endline) at (12.3,0)
    {Endline survey\\2--14 days later};
  \draw[arrow] (census) -- (outreach);
  \draw[arrow] (outreach) -- (campaign);
  \draw[arrow] (campaign) -- (endline);
  \node[font=\footnotesize, text=gray!70!black, below=3pt of campaign]
    {Two 12-day waves: Oct.--Nov. 2016};
\end{tikzpicture}
\caption{Experiment Timeline and Data Collection}


\label{fig:timeline}
\end{figure}

Our analysis combines administrative take-up data at treatment sites with surveys collected before and after the campaign. Appendix \ref{app:sample-construction} and Figure \ref{fig:sample-consort} provide additional details on how each sample was constructed and how they correspond to the surveys.


\begin{enumerate} 
\item \emph{Household census.}
We conducted a census of all adults, aged 18 or older, in the 144 selected study communities. Enumerators visited each household to record geographic coordinates and collect data on basic demographics and phone ownership. We use the census data for random assignment of the SMS treatment, and to draw subsamples for the baseline survey and monitoring of deworming take-up. We further use the data to check for balance of individual and community characteristics across incentive and distance conditions and for control variables in the main estimation. 

\item \emph{Baseline survey.} We surveyed a random sample of adults to measure prior deworming experience, knowledge about treatment, and social norms surrounding treatment. This data are used to provide descriptives of the context and for additional balance checks. 

\item \emph{Administrative take-up data.} We directly monitored deworming take-up at PoTs, avoiding relying on individual self-reports for our primary outcome.\footnote{This is critical for that demand effects in survey responses cannot mechanically generate the main reduced form take-up results.} Enumerators used tablets pre-loaded with the census roster to confirm sampled adults' identity at the PoT and record whether they received deworming treatment. The monitoring sample consists of baseline individuals, SMS treatment and control individuals and an additional sample of non-SMS individuals, including phone and non-phone owners. This data are used to estimate main treatment effects on deworming take-up and to assess the correctness of adults' beliefs about peers' deworming take-up.


  \item \emph{Endline survey.} We conducted an in-depth survey with a random subsample of adults whose take-up was monitored 2–14 days after the deworming campaign. It contains three key modules:
  \begin{itemize} 
      \item \emph{Observability module.} Half of respondents are
  randomly assigned to an observability module eliciting beliefs
  about peers' take-up decisions.\footnote{The remaining half are assigned to a perceived observability module, exploring respondents' beliefs about what others know about their own deworming status.} Respondents were shown a random subset of ten adults from their community and, for each person they recognized, asked: ``Do you think this person came for deworming?''. The module is designed to capture how much respondents know about others’ deworming behavior within their community, reflecting the extent to which such actions are visible.
        
      \item \emph{Gift choice and WTP modules.} All respondents
  chose between a bracelet and a calendar as a thank-you gift. In
  control communities we augmented this with a willingness-to-pay exercise, offering random cash incentives to switch items. Together these data test non-social-image mechanisms and pin down the private utility difference between incentives in the structural model.                                        
  \end{itemize}
  Endline data are also used to verify correct implementation of
  the incentive treatments and to measure retention of the
  incentive items.    

\end{enumerate}

\subsection{Sample Definition and Randomization Checks} 

The analysis sample that measures our primary outcome, individuals' deworming status, comprises 9,805 adults who did not receive an SMS treatment and for which we monitored their take-up directly at the PoT (``Take-up analysis sample''). For observability and belief outcomes this is 1,141 individuals in the endline survey that receive the observability module but do not receive an SMS treatment (``Endline observability sample''). Table~\ref{tab:stacked-sample-balance-table} reports control group means and standard errors for the two main outcome samples across each baseline covariate, together with pairwise differences between treatment and control groups and the associated $p$‑values, presented separately for the Far and Close conditions. The final column shows $F$‑tests of joint significance across all treatments.  All $p$‑values exceed 0.10 except for two covariates: (i) Distance to PoT in both samples; and (ii) Female, driven by fewer women in the Far-Bracelet communities in the endline observability sample. 

Focusing on the first imbalance, within the Far condition, communities assigned to Ink are on average 344 meters farther from the PoT than Far–Control communities ($p = 0.02$). To address this, we include continuous community‑to‑PoT distance as a control in robustness specifications of the reduced form analysis; the structural model conditions directly on continuous distance. The second imbalance only occurs in the smaller endline observability sample and is controlled for by a gender dummy included in the post-double LASSO selected set of controls in both take-up and observability regressions.

Overall, sample characteristics are well balanced. Out of 80 comparisons in Table \ref{tab:stacked-sample-balance-table}, 8 differ at the 10 percent level and 5 at the 5 percent level---in line with the 8 and 4 differences expected by chance. Additionally, because our identification strategy exploits variation both within distance conditions (across incentive arms) and within incentive arms (across distance), Table \ref{tab:stacked-sample-dist-balance-table} reports balance across distance within each incentive arm.  Of these 32 additional comparisons, only one differs at the 10 percent level and none at the 5 percent level, again consistent with chance.

\subsubsection{Sample Attrition}
\label{sec:sample-attrit}
We observe two instances of sample attrition. First, in the monitored deworming sample, 989 individuals assigned to be monitored were not, due to a SurveyCTO error.\footnote{EXPLAIN ERROR} Second, amongst the 3,729 individuals assigned to receive the endline observability module, 255 erroneously did not receive the component. Table~\ref{tab:takeup-attrit} examines whether treatment assignment predicts attrition from the monitored sample (shown in Figure \ref{fig:sample-consort}). The dependent variable is an indicator equal to one if an observation was dropped from the sample. Across all treatment arms---Bracelet, Calendar, and Ink---the estimated coefficients are small in magnitude and statistically indistinguishable from zero, both in the combined specification and when the sample is split by distance. Taken together, these results provide no evidence that treatment assignment induced differential attrition, supporting the integrity of the experimental sample.

Table~\ref{tab:endline-obs-attrit} shows that attrition into the endline observability sample is more complex. Bracelet-arm respondents are significantly less likely to be missing from the observability survey module, with a pooled gap of 6.1 percentage points and a close-distance gap of 8.9 percentage points. This differential is primarily mechanical and pre-determined: the SurveyCTO knowledge list was pre-loaded only for primary endline respondents, so backup households---surveyed when a primary respondent could not be located---had this section silently skipped. Control households had the highest backup rates $(\approx19\%)$ compared to bracelet households $(\approx13\%)$, plausibly because bracelet participants were more engaged and easier to locate. Controlling for backup status absorbs roughly half the treatment gap, and crucially, own deworming status does not predict selection into the sample (coefficient = 0.011, $p = 0.37$). A residual gap of approximately 4.5 percentage points remains after conditioning on backup status, though this is unlikely to reflect endogenous selection on deworming behavior. To account for this residual imbalance, we present Lee bounds on the observability outcomes in the appendix.

\subsection{Implementation Fidelity and Adoption of Signals} %UNDERSTANDING OF INCENTIVES -- TO WHAT EXTENT want to presuppose SIGNAL HERE instead of INCENTIVES



\label{tab:incentive_check_smscontrol}
Endline survey responses confirm that the deworming campaign was implemented as intended
(Table \ref{tab:implement-bal}). Across both Close and Far distance groups, treatment arms
show no statistically significant differences from the control in CHV visitation rates
(joint $p$-values of 0.809 and 0.613, respectively), with an overall sample mean of 82
percent. Rates of community announcement, CHV sharing of deworming practices, and CHV
communication of treatment locations are similarly balanced across arms, with all joint
tests yielding $p > 0.28$.


Similarly, table \ref{tab:incentive-check} demonstrates that the three incentive items were distributed correctly at the point of treatment (PoT): more than 95 percent 
of respondents in each respective arm reported receiving their assigned item. Recognition of 
item meaning was similarly high, with 94–97 percent of respondents correctly associating the 
ink and bracelet with deworming participation, and 89 percent identifying the calendar 
as a signal of treatment attendance.

Given that indelible ink fades within approximately three days on skin and two weeks on nails, 
only 14 percent of respondents displayed a visible ink mark at endline. Retention of the 
remaining incentives was substantially higher: 79 percent of respondents still possessed the 
bracelet and 87 percent retained the calendar. Among bracelet recipients, 46 percent were still 
wearing it on their wrist at endline, and 74\% of calendar recipients had it displayed on their 
wall (Table \ref{tab:incentive_check_smscontrol}).

The high rates of bracelet retention and active wearing suggest that the incentive was 
positively valued by recipients, rather than carrying negative utility through embarrassment 
or social stigma. The bracelet was also the most publicly visible incentive: 95 percent of 
endline respondents—irrespective of own deworming status—reported observing it within their 
community. The calendar, by contrast, exhibited high retention but limited public visibility, 
making it an appropriate active-control comparison for isolating any self-signaling 
(self-image) effects attributable to the bracelet.


% APPENDIX STARTS HERE

\subsection{Sample Construction}
\label{app:sample-construction}
\begin{landscape}
    \begin{figure}
      \centering                                                   
      \adjustbox{max width=\linewidth}{\begin{tikzpicture}[
  node distance = 7mm and 16mm,
  font          = \small,
  %
  % ── Node styles ─────────────────────────────────────────────────
  box/.style      = {draw, rounded corners=2pt, align=center,
                     text width=5.2cm, inner sep=5pt},
  analysis/.style = {box, fill=analysisgreen},
  mon/.style      = {box, fill=monblue},
  issue/.style    = {box, fill=issueorange},
  note/.style     = {box, dashed},
  arrow/.style    = {-Latex, line width=0.7pt},
  darr/.style     = {-Latex, line width=0.7pt, dashed},
  colgroup/.style = {draw, rounded corners=3pt, inner sep=10pt},
]

%% ═══════════════════════════════════════════════════════════════════
%% SHARED TOP — Census frame
%% ═══════════════════════════════════════════════════════════════════

\node[draw, rounded corners=2pt, align=center,
      text width=23cm, inner sep=6pt]
  (census)
  {\textbf{Household census:} 144 clusters,\ 39{,}301 adults (18+)};

%% ═══════════════════════════════════════════════════════════════════
%% COLUMN A — Baseline
%%   xshift = -6.5 cm keeps this column on the left
%% ═══════════════════════════════════════════════════════════════════

\node[box, below=5mm of census, anchor=north, xshift=-7.5cm]
  (bl-sel)
  {Baseline roster: 4{,}185};

\node[box, below=of bl-sel] (bl-tgt)
  {Selected: 2{,}160};

\node[box, below=of bl-tgt] (bl-reach)
  {Reached: 2{,}062};

\node[analysis, below=of bl-reach] (bl-comp)
  {\textbf{``Baseline sample''}\\
  Completed baseline: 1{,}995};

%% Side branch: not completed (exits left from Reached)
\node[box, left=16mm of bl-reach, anchor=east] (bl-nc)
  {Not completed: 67\\(refused 29; unable 38)};

%% ── Arrows ───────────────────────────────────────────────────────
\draw[arrow] (census.south -| bl-sel.north) -- (bl-sel.north);
\draw[arrow] (bl-sel)   -- (bl-tgt);
\draw[arrow] (bl-tgt)   -- (bl-reach);
\draw[arrow] (bl-reach) -- (bl-comp);
\draw[arrow] (bl-reach.west) -- (bl-nc.east);

%% ── Column A group box ───────────────────────────────────────────
\node[colgroup,
      fit=(bl-tgt)(bl-reach)(bl-comp)(bl-nc),
      label={[font=\small\bfseries, xshift=-1cm]above:Baseline survey}] {};

%% ═══════════════════════════════════════════════════════════════════
%% COLUMN B — Monitored sample
%%   xshift = 0 cm (center)
%% ═══════════════════════════════════════════════════════════════════

% Merged totals (fed from bl-sel in A and el-frame in C)
\node[mon, below=40mm of census, anchor=north] (mon-uniq)
  {Total unique individuals: 12{,}827};

\node[analysis, below=of mon-uniq] (mon-ana)
  {\textbf{``Take-up analysis sample''}\\
  Excl.\ SMS treatment: 9{,}805};

%% ── Arrows ───────────────────────────────────────────────────────
% Baseline roster enters from col A
\draw[arrow] (bl-sel.east) -- (mon-uniq.north);

\draw[arrow] (mon-uniq) -- (mon-ana);

%% ── Column B group box ───────────────────────────────────────────
\node[colgroup,
      fit=(mon-uniq)(mon-ana),
      label={[font=\small\bfseries, yshift=13mm]above:Administrative take-up data}] {};

%% ═══════════════════════════════════════════════════════════════════
%% COLUMN C — Endline
%%   xshift = 6.5 cm keeps this column on the right
%% ═══════════════════════════════════════════════════════════════════

\node[box, below=5mm of census, anchor=north, xshift=11cm]
  (el-frame)
  {Endline roster: 11{,}166
  };

\node[box, below=of el-frame] (el-sel)
  {Selected: 4{,}220};

\node[box, below=of el-sel] (el-reach)
  {Reached: 3{,}830};

\node[box, below=of el-reach] (el-comp)
  {Completed endline: 3{,}729};

%% PAP deviation note (exits left from Completed)
\node[note, left=8mm of el-comp, anchor=east, text width=3.5cm] (el-pap-dev)
  {\textbf{PAP deviation:}\\
  3{,}600 completions (25/cluster $\times$ 144)\\
  {\scriptsize Overshoot: $\sim$2 backups/cluster surveyed}};

\node[analysis, below=of el-comp] (el-comp-nosms)
  {\textbf{``Endline sample''}\\
  Excl.\ SMS treatment: 2{,}659};

%% Side branch: not completed (exits right from Reached)
\node[box, right=16mm of el-reach, anchor=west] (el-nc)
  {Not completed: 101\\(refused 23; unable 78)};

%% WTP sample (below endline sample, control only)
\node[analysis, below=of el-comp-nosms] (wtp)
  {\textbf{``WTP sample''}\\
  Control clusters only (24 of 144)\\
  Targeted: 600, Reached: 474 \\
  Completed: 467};

%% ── Arrow ────────────────────────────────────────────────────────
\draw[arrow] (el-comp-nosms) -- (wtp);

%% Side branch: missing knowledge module
\node[issue, right=16mm of el-comp, anchor=west] (el-miss)
  {Missing knowledge\\module: 255};

%% Completed knowledge module (complement of el-miss)
\node[box, below=of el-miss] (el-kmod)
  {Observability module: 1{,}784 \\
  (Randomly assigned)
  };

%% Observability analysis sample
\node[analysis, below=of el-kmod] (el-know)
  {\textbf{``Endline observability sample''}\\
  Excl.\ SMS treatment: 1{,}141};


%% ── Arrows ───────────────────────────────────────────────────────
\draw[arrow] (census.south -| el-frame.north) -- (el-frame.north);
\draw[arrow] (el-frame) -- (el-sel);
\draw[arrow] (el-sel)   -- (el-reach);
\draw[arrow] (el-reach) -- (el-comp);
\draw[arrow] (el-reach.east) -- (el-nc.west);
\draw[darr] (el-comp.west) -- (el-pap-dev.east);
\draw[arrow] (el-comp)  -- (el-comp-nosms);
\draw[arrow] (el-comp.east) -- (el-miss.west);
\draw[arrow] (el-comp.east) -- ++(4mm, 0) |- (el-kmod.west);
\draw[arrow] (el-kmod) -- (el-know);

%% ── Column C group box ───────────────────────────────────────────
\node[colgroup,
      fit=(el-sel)(el-reach)(el-comp)(el-comp-nosms)(wtp)(el-miss)(el-nc)(el-kmod)(el-know),
      label={[font=\small\bfseries, xshift=1cm]above:Endline survey}] {};

%% ── Cross-column: el-frame forks → Total unique and → Dropped ──
% Split point midway between cols B and C, at el-frame height
\coordinate (el-split) at ($(mon-uniq.north |- el-frame.west) + (5.5cm, 0)$);

% Dropped box hangs below the split point (between cols B and C)
\node[issue, anchor=north,
      text width=2.2cm, inner sep=3pt] (mon-drop)
  at ($(el-split) + (0, -12mm)$)
  {Dropped\\(SurveyCTO error): 989};

% Stem: el-frame → split (no arrowhead)
\draw[line width=0.7pt] (el-frame.west) -- (el-split);
% Branch 1: left-then-down to Total unique
\draw[arrow] (el-split) -- (mon-uniq.north);
% Branch 2: straight down to Dropped
\draw[arrow] (el-split) -- (mon-drop.north);

\end{tikzpicture}
}
      \caption{Sample Construction Diagram}    
      \label{fig:sample-consort}
    \end{figure}  
  \end{landscape}


\clearpage
\newpage

% ── Baseline sample ───────────────────────────────────────────────
Figure \ref{fig:sample-consort} displays how each survey maps onto our main analysis samples:

\emph{Baseline survey.} The baseline sampling frame was drawn from the household
census roster of 39{,}301 adults across 144 clusters.  From this frame, 4{,}185
individuals were selected by randomly drawing 15 households per cluster (additionally, we drew two backup 
households per cluster).  Enumerators attempted to reach a
target of 2{,}160 individuals; of the 2{,}062 successfully located, 38 were
unable to participate (e.g. due to illness or language barriers) and 29
refused consent, yielding 1{,}995 completed baseline interviews that
constitute the ``Baseline sample''.

The full 4{,}185-person baseline roster plays a second role beyond the baseline
survey itself: it was added to the administrative monitoring lists used to
record deworming take-up at points of treatment on the day of the campaign.  As
a result, the baseline roster contributes approximately one-fifth of the broader
monitored take-up sample, discussed in the next section.

% ── Endline sample ────────────────────────────────────────────────

\emph{Endline survey.} The endline sampling roster consisted of 11{,}166
individuals drawn from the household census roster.  From this roster, 4{,}220
individuals were selected, comprising a main roster of approximately 25 per
cluster plus roughly two backup individuals per cluster, and enumerators
successfully located 3{,}830 of them.  Of those reached, 78 were unable to
participate and 23 refused consent, yielding 3{,}729 completed endline
interviews.  The PAP pre-specified 3{,}600 completions (25 per cluster $\times$
144 clusters), the overshoot is driven by the (mistaken) systematic surveying of
all backup individuals. Excluding the SMS treatment arm, which is set aside for
the SMS analysis, the ``Endline sample'' used for the main analysis comprises
2{,}659 respondents.  The full 11{,}166-person endline roster is added to the
administrative monitoring lists for their deworming status to be recorded at
points of treatment during the campaign.


The endline survey generates three further subsamples used throughout the
analysis.  First, a random half of endline respondents were assigned to an
observability module eliciting beliefs about peers' deworming take-up (1{,}784
respondents), of whom 1{,}141 remain after excluding the SMS treatment arm,
the ``Endline observability sample''. There are 255 individuals with
missing observability module data due to a programming error, discussed in Section \ref{sec:sample-attrit}.  Second, all endline respondents stated a preference between a
bracelet and a calendar as a thank-you gift (\emph{gift choice}).  Our analysis 
using this gift-choice data focuses on individuals that did not choose to deworm
in the control group (1{,}350 individuals). Third, we conducted a
willingness-to-pay (WTP) exercise for the bracelet gift in the control clusters,
which is used to estimate a WTP model for the bracelet.  In this WTP exercise,
we attempted to reach 600 individuals across 34 control clusters, of whom we
successfully reached 474 and obtained completed WTP data from 467, ``WTP sample''.


Because the endline roster is drawn randomly from the household census, endline
respondents partly overlap with the baseline sample, but we are unable to link
baseline and endline records at the household level.  The two surveys therefore
provide independent cross-sectional snapshots of the same communities rather
than a panel.

% ── Take-up analysis sample ───────────────────────────────────────

\emph{Take-up analysis sample.}
Deworming take-up was recorded administratively at points of treatment for all
individuals on the monitoring lists, which were assembled from the baseline
roster (4{,}185 individuals) and the endline roster (11{,}166 individuals). Unfortunately, 
a SurveyCTO coding error caused 989 individuals from the endline roster to be omitted from the 
administrative monitoring lists, and thus they are absent from the take-up analysis sample, further discussion including differential attrition tests can be found in Section \ref{sec:sample-attrit}.
Because the two rosters were drawn independently from the same census, some
individuals appear on both lists; after deduplication the combined monitoring
universe comprises 12{,}827 unique individuals.
Excluding the SMS treatment arm, which was assigned to receive social
information or reminder messages and is reserved for the SMS analysis, yields
the take-up analysis sample of 9{,}805 individuals used throughout the
main results.

  
